% 第三部：攻擊的標的（第壹部分，1353–1359）。出處：課程倉 W06_1019_導讀問答.md「強式與弱式」一則（已查核）。
\chapter{攻擊的標的：違憲審查的界定}

\qaentry{第三部之一}{強式與弱式的區分標準}

\begin{qaquestion}
本文攻擊的是哪一種違憲審查？強式與弱式以什麼為區分標準？
\end{qaquestion}

第壹部分開頭先限定範圍：本文處理法院對立法的審查，對行政行為的司法審查依據的是法治原則，不在本文範圍\pg{1353--1354}。在對立法的審查中，作者認為最重要的區分是強式與弱式\pg{1354}。強式的定義是：

\begin{quote}
In a system of strong judicial review, courts have the authority to decline to apply a statute in a particular case (even though the statute on its own terms plainly applies in that case) or to modify the effect of a statute to make its application conform with individual rights (in ways that the statute itself does not envisage).\pg{1354}
\end{quote}

法院另得經由遵循先例原則與爭點效，使該法律實際上成為具文。弱式的定義是：

\begin{quote}
In a system of weak judicial review, by contrast, courts may scrutinize legislation for its conformity to individual rights but they may not decline to apply it (or moderate its application) simply because rights would otherwise be violated.\pg{1355}
\end{quote}

區分的標準因此是法院在個案中得否拒絕適用法律，或修正其效果。註 22 另指出，強弱之分與司法至上（judicial supremacy）屬不同問題\pg{1354}。

\qaentry{第三部之二}{各國制度的歸類}

\begin{qaquestion}
本文如何把美國、加拿大、英國、紐西蘭與部分歐洲法院排在強弱之間？加拿大有儘管條款，為何仍算強式？
\end{qaquestion}

\begin{qafigure}{違憲審查的五種形式與原文所舉之例（1354--1357）。\label{fig:spectrum}}
% 違憲審查的五種形式，由強至弱（1354–1357）。灰階；本文攻擊的標的以強調色括號標出。
\begin{tikzpicture}[qa]
  \foreach \i/\n/\d/\e in {
      0/更強/{將法律自法典中剔除}/{部分歐洲法院\\1354--1355 註 23},
      1/強式/{個案中拒絕適用或修正其效果；經先例使之成為具文}/{美國\\1355},
      2/居間，計為強式/{拒絕適用；立法機關得以儘管條款排除}/{加拿大憲章第 33 條\\1356--1357},
      3/弱式/{不得拒絕適用；得為不相容宣告}/{英國 1998 年人權法\\1355},
      4/更弱/{只得作合於權利的解釋}/{紐西蘭 1990 年權利法案法\\1356}}{
    \node[nbox,text width=25mm,minimum height=21mm,anchor=north west] (b\i) at (\i*2.95,0) {\n\\[2pt]{\scriptsize \d}};
    \node[lbl,text width=26mm,align=left,anchor=north west] at (\i*2.95,-2.35) {\e};
  }
  \draw[figGray,line width=0.6pt,-{Stealth[length=4.5pt]}] (0,0.45) -- node[lbl,above] {法院的權限由強至弱} (14.6,0.45);
  \draw[figE,line width=0.9pt] (0.05,-3.25) -- (8.5,-3.25);
  \node[font=\scriptsize,text=figE,anchor=north west] at (0.05,-3.32) {本文攻擊的標的};
\end{tikzpicture}

\end{qafigure}

美國法院得拒絕適用，依原文似不能將法律剔除，但實際效果與剔除相去不遠\pg{1355}。註 24 討論美國是否屬於更強的形式，稱此事並無明確答案：Dickerson v.\ United States 案 Scalia 大法官不同意見書的末句，似乎顯示被最高法院認定違憲的法律仍可供法院援用；McCorvey v.\ Hill 案中第五巡迴上訴法院認定德州墮胎法已遭默示廢止，作者細讀後指出，默示廢止出於 Roe 之後德州的墮胎立法，並非 Roe 判決本身\pg{1355}。

加拿大權利與自由憲章第 33 條使立法機關得明示宣告法律不受憲章第 2 條或第 7 至 15 條的拘束。作者仍將加拿大計為強式，理由是該條款實際上甚少援用\pg{1356--1357}。註 34 回應 Goldsworthy「該條款足以回應民主疑慮」之說：援用該條款等於要立法機關表示自己不認為憲章權利重要，而法院與立法機關的對立，通常是雙方都重視該權利而對其內容理解不同\pg{1357}。

\qaentry{第三部之三}{作者接受的違憲審查}

\begin{qaquestion}
作者是否反對一切形式的違憲審查？
\end{qaquestion}

第參部分前段，作者表示尊重一項論證：議員未必看得出法案所涉的權利問題，也未必預見法律適用時會發生何種權利問題，因此宜有一套機制，讓公民在問題出現時提請眾人注意\pg{1370}。他接著寫：

\begin{quote}
But this is an argument for weak judicial review only, not for a strong form of the practice in which the abstract question of right that has been identified is settled in the way that a court deems appropriate.\pg{1370}
\end{quote}

這項論證所支持的，是英國那種由法院宣告存在重要權利問題的制度，或由檢察總長審查法律案並公開指出其中權利問題的更弱安排；註 63 引紐西蘭權利法案法第 7 條\pg{1370}。

\qaentry{第三部之四}{另三項區分}

\begin{qaquestion}
除強弱之外，第壹部分還作了哪些區分？作者如何取捨？
\end{qaquestion}

\begin{longtable}{@{}L{0.30\linewidth}L{0.62\linewidth}@{}}
\toprule
區分 & 作者的取捨 \\
\midrule
\endhead
以權利為對象的審查與一般的憲法審查（1357--1358） & 本文處理前者。聯邦主義事項由大法官投票決定，作者另表保留；註 35 指出最著名的違憲審查判決 Marbury v.\ Madison 與個人權利無關 \\
事後的個案審查與事前的抽象審查（1358--1359） & 作者表明專注於事後審查，對事前審查不再多述 \\
普通法院與專設的憲法法院（1359） & 以 Goodridge 為普通法院行使審查之例，對專設憲法法院只寫一段 \\
\bottomrule
\end{longtable}

\begin{readernote}
後兩項區分與中華民國的制度直接相關：憲法法庭是專設的憲法法院，國家最高機關與立法委員現有總額四分之一以上聲請的法規範憲法審查（憲法訴訟法第 47 條、第 49 條）不以具體案件為前提。本文第肆部分 A 節批評的「個案取向」，在這類審查中本不存在；可以請聽者評估，本文其餘的論證對這類制度的適用程度。
\end{readernote}
